\chapter{آلة بلا قائد أوركسترا}
% full version: chapters/02_glutcomputer.tex

في حاسوبك قائد أوركسترا: مولّد الساعة. مليارات المرات في الثانية يعطي الإشارة، فتتبدل كل عناصر الدائرة معًا، بخطوة واحدة. أما في الدماغ فلا قائد. كل عصبون يعمل وحده: تصله الإشارات متى اتفق، وتخرج منه متى اتفق. لنرَ ما الذي يمكن حسابه في آلة كهذه إذا أخذنا عصبونات \nt{GLUT} وحدها، أي ”الزوائد“ وحدها.

\section*{الدلو المثقوب}

من المريح أن نتخيل العصبون دلوًا مثقوبًا. كل نبضة تصل هي مغرفة ماء. والماء يتسرب شيئًا فشيئًا من الثقب. إذا توالت المغارف بسرعة، ارتفع الماء إلى الحافة، فينقلب الدلو: ينقر العصبون، ويفرغ الدلو، ويبدأ كل شيء من جديد. وإذا جاءت المغارف متباعدة، تسرب الماء ولم يحدث شيء.

\pencilfig{2}{\pencilart{2}}{العصبون دلو مثقوب: قطرات المداخل تملؤه والماء يتسرب؛ إذا بلغ الخط، كانت نقرة.}

الثقب هو الفارق الأساسي بين العصبون والبوابة المنطقية. العصبون لا يتذكر النبضة الواصلة إلى الأبد، بل بضعة أجزاء من ألف من الثانية. كل ما يصل ضمن هذه النافذة يُجمع؛ وما سبقها يُنسى.

\section*{”أو“ و”و“}

إذا كفت مغرفة واحدة لقلب الدلو، استجاب العصبون لأي مدخل. هذه ”أو“. وإذا لزمت مغرفتان، بدأ الأمر المثير: سؤال ”هل وصلت الاثنتان؟“ بلا قائد لا معنى له ما لم نحدد \emph{خلال أي مدة}. لن ينقر العصبون إلا إذا وصلت النبضتان في وقت واحد تقريبًا، قبل أن تتسرب المغرفة الأولى. هذه ”و في الزمن“: كاشف التزامن. عند بعض العصبونات تكون النافذة بضعة أجزاء من ألف من الثانية، وعند عصبونات السمع أجزاءً من جزء من ألف.

\section*{ما لا تقدر عليه ”الزوائد“ وحدها}

حاول أن تبني ”لا“: عصبونًا يعمل حين يصمت المدخل. لن تنجح: الصمت لا يصب ماء في الدلو. وهذا صحيح لا لعصبون واحد فحسب، بل لأي شبكة من ”الزوائد“ وحدها: إذا قويت المداخل، لن يعمل أي عصبون في الشبكة أقل مما كان. ومن هنا أسوأ خصائصها: \emph{لا يمكن إطفاء النشاط بالنشاط}. الحريق المشتعل لا يوقفه إلا شيء واحد: الكف عن رمي الحطب.

وجدت الطبيعة طريقًا التفافيًا. في كثير من أعضاء الحس تنقل مجموعتان من العصبونات المنبه نفسه. في العين تُبلغ بعض الخلايا ”صار أكثر إضاءة“، وأخرى ”صار أكثر عتمة“. فإذا كان لديك بدل السلك الواحد زوج سلكين ”نعم“ و”لا“، جاءت ”لا“ مجانًا: يكفي أن تبدل السلكين. وبأزواج كهذه يمكن تجميع أي دائرة منطقية من ”الزوائد“ وحدها، بل ومعرفة أن الحساب انتهى: في كل زوج اشتعل سلك واحد بالضبط. هكذا يبني المهندسون الدوائر غير المتزامنة التي تعمل صحيحًا مع أي تأخيرات.

\section*{الزمن من التأخيرات}

بلا قائد لا توجد ساعة مشتركة، لكن يمكن حساب الزمن مع ذلك. الصوت الآتي من الجانب يصل إلى الأذن القريبة قبل البعيدة بأجزاء من جزء من ألف من الثانية. لنمدّ من الأذن اليسرى سلكًا من اليسار إلى اليمين على طول صف من كواشف التزامن، ومن اليمنى سلكًا من اليمين إلى اليسار. ستجري الإشارتان إحداهما نحو الأخرى، وتلتقيان في النقطة التي تبلغها كلتاهما معًا. هناك بالضبط يعمل الكاشف. تحوّل فرق الزمن إلى \emph{رقم} العصبون الذي عمل، أي إلى اتجاه الصوت. وتبلغ الدقة عشرات من أجزاء المليون من الثانية، مع أن كل عصبون بمفرده أقل دقة بكثير: ما ينقذ الموقف هو كثرة الكواشف.

\pencilfig{102}{%
% delay line: the left-ear wire runs right along the top, the right-ear wire runs left along the bottom
\draw[pen=pcBlue] (0,0.6) -- (5.6,0.6); \draw[pen=pcBlue] (6.0,-0.6) -- (0.4,-0.6);
\foreach \i in {1,...,7}{
  \pgfmathsetmacro{\x}{0.75*\i}
  \draw[pen=pcBlue] (\x,0.6) -- (\x,0.24); \draw[pen=pcBlue] (\x,-0.6) -- (\x,-0.24);
  \ifnum\i=5 \draw[dotpen=pcAmber, wash=pcAmber, line width=1.1pt] (\x,0) circle (0.2);
  \else \draw[dotpen=pcBlue, wash=pcBlue] (\x,0) circle (0.2); \fi}
\draw[arr=pcBlue] (0.95,0.78) -- (1.75,0.78); \draw[arr=pcBlue] (5.3,-0.78) -- (4.5,-0.78);
\node[lbl, anchor=east] at (-0.05,0.6) {الأذن اليسرى};
\node[lbl, anchor=west] at (6.05,-0.6) {الأذن اليمنى};
\node[lbl, text=pcAmber!70!black, anchor=south] at (3.75,0.95) {هنا التقتا};
\draw[arr=pcAmber!70!black] (3.75,0.95) -- (3.75,0.68);
% sound source on the left
\draw[penlite] (-1.25,-0.25) -- (-1.05,-0.25) -- (-0.8,-0.45) -- (-0.8,0.05) -- (-1.05,-0.15) -- (-1.25,-0.15) -- cycle;
\foreach \r in {0.2,0.35}{\draw[penlite] ($(-0.75,-0.2)+(-40:\r)$) arc (-40:40:\r);}
\node[lbl, anchor=north] at (-0.95,-0.5) {صوت من اليسار};
}{الصوت الآتي من اليسار يصل إلى الأذن اليسرى أولًا، فتقطع إشارتها مسافة أطول، ويقع اللقاء يمين المنتصف. رقم العصبون الذي عمل هو اتجاه الصوت.}

\section*{الذاكرة نار مخيّم}

لنأخذ مجموعة من العصبونات يثير بعضها بعضًا بقوة. إذا أشعلناها بإشارة قصيرة، غذّت نفسها بنفسها، كنار تلقي في ذاتها وقودها. انتهت الإشارة، والمجموعة لا تزال تشتعل: إنها تتذكر أن الإشارة كانت. هذه خلية ذاكرة. وإذا كانت الإشارة أقصر مما ينبغي، لم تلحق النار أن تتقد فتنطفئ.

لكن هنا المشكلة: كيف نمحو ذاكرة كهذه؟ بالنشاط لا سبيل إلى ذلك. لا يبقى إلا انتظار أن تتعب العصبونات وتخبو النار.

\section*{مؤقّت يطلقه حدث}

لنصل المجموعات في سلسلة: الأولى تشعل الثانية، والثانية الثالثة. ندفع الأولى، فتجري الموجة في السلسلة كما في أحجار الدومينو. ورقم الحجر الساقط يخبر كم مر من الوقت منذ الدفعة. ليست هذه ساعة تدق دائمًا، بل مؤقت يشغّله حدث ويعمل مرة واحدة. عند الطيور المغردة، أثناء الغناء، تعمل عصبونات إحدى المناطق واحدًا تلو الآخر بترتيب صارم، وهذا يشبه كثيرًا سلسلة كهذه.

إذن تصنع ”الزوائد“ وحدها أشياء كثيرة. وينقصها ثلاثة: اختيار واحد من كثير، ومحو الذاكرة بأمر، ومنع الشبكة من الاحتراق. والثلاثة تأتي بها ”النواقص“، وهي موضوع الفصل التالي.

\fullversion{2}
